% fake_monster_sibling_costello.tex
% Standalone note on the Fake-Monster denominator, the K3 x E
% Delta_5 branch, and the conditional dimension-five comparison.

\documentclass[11pt]{amsart}
\usepackage[localtheorems]{raeez-math-template}

\newtheorem{theorem}{Theorem}[section]
\newtheorem{proposition}[theorem]{Proposition}
\newtheorem{lemma}[theorem]{Lemma}
\newtheorem{corollary}[theorem]{Corollary}
\theoremstyle{definition}
\newtheorem{definition}[theorem]{Definition}
\newtheorem{conjecture}[theorem]{Conjecture}
\theoremstyle{remark}
\newtheorem{remark}[theorem]{Remark}

\providecommand{\C}{\mathbb{C}}
\providecommand{\Z}{\mathbb{Z}}
\providecommand{\fg}{\mathfrak{g}}
\providecommand{\cO}{\mathcal{O}}
\providecommand{\Leech}{\Lambda_{24}}
\providecommand{\IIlor}{\mathrm{II}_{25,1}}
\providecommand{\IIpar}{\mathrm{II}_{26,2}}
\providecommand{\PhiFM}{\Phi_{12}}
\providecommand{\gFM}{\mathfrak{g}_{\mathrm{FM}}}
\providecommand{\gDelta}{\mathfrak{g}_{\Delta_5}}
\providecommand{\PhiFA}{\Phi^{\mathrm{FA}}}
\providecommand{\SpCh}{\mathrm{Sp}^{\mathrm{ch}}}
\providecommand{\Perf}{\mathrm{Perf}}
\providecommand{\CoHA}{\mathrm{CoHA}}
\providecommand{\kBKM}{\kappa_{\mathrm{BKM}}}
\providecommand{\kch}{\kappa_{\mathrm{ch}}}
\providecommand{\kcat}{\kappa_{\mathrm{cat}}}
\providecommand{\kfib}{\kappa_{\mathrm{fiber}}}
\providecommand{\kHeis}{\kappa_{\mathrm{ch}}^{\mathrm{Heis}}}
\providecommand{\rk}{\mathrm{rk}}
\providecommand{\wt}{\mathrm{wt}}

\begin{document}

\title{Fake-Monster Denominator Data and the Dimension-Five Comparison}
\author{}
\date{}
\maketitle

\begin{abstract}
The Fake-Monster Lie algebra of Borcherds lives on
\(\IIlor=\Leech\oplus U\), has Weyl vector at the Leech cusp, and has
denominator \(\PhiFM\) of weight \(12\) on the type-IV domain attached to
\(\IIpar=\IIlor\oplus U\). This datum is separate from the K3
\(\Delta_5\) denominator governing the \(K3\times E\) BKM branch. The
single-K3-fibre route is excluded by rank and signature. A
\(K3_1\times K3_2\times E\) comparison is a conditional dimension-five
datum requiring a primitive Leech-cusp embedding, a Hall comparison, and a
factorisation-algebra specialisation whose denominator is \(\PhiFM\).
\end{abstract}

\section{Borcherds Data}

Let \(\Leech\) be the positive-definite Leech lattice of rank \(24\), and
let \(U=\Z e\oplus \Z f\) be the hyperbolic plane with
\[
 e^2=f^2=0,\qquad (e,f)=1 .
\]
Set
\[
 \IIlor=\Leech\oplus U,\qquad \rk \IIlor=26,\qquad
 \operatorname{sign}(\IIlor)=(25,1),
\]
with the sign convention used in Borcherds' Fake-Monster construction.
The primitive isotropic vector
\[
 \rho=(0;1,0)\in \Leech\oplus U
\]
is the Weyl vector at the Leech cusp: \(\rho^2=0\), and
\(\rho^\perp/\Z\rho\simeq \Leech\). The real simple roots are the
norm-\(2\) Leech-cusp roots. The imaginary simple roots are positive
multiples of \(\rho\), with multiplicities read from
\[
 \eta(\tau)^{-24}
 =
 q^{-1}\prod_{n\geq1}(1-q^n)^{-24}
 =
 q^{-1}+24+324q+3200q^2+\cdots .
\]
Write
\[
 \eta(\tau)^{-24}=\sum_{m\geq -1} c(m)q^m .
\]

\begin{theorem}[Borcherds denominator for the Fake Monster]
\label{thm:fm-borcherds-denominator-data}
On the tube domain of the Leech cusp in
\(\IIlor\otimes\C\), the Fake-Monster denominator is
\[
 \PhiFM(Z)
 =
 e^{2\pi i(\rho,Z)}
 \prod_{\substack{\alpha\in\IIlor\\(\alpha,\rho)>0}}
 \left(1-e^{2\pi i(\alpha,Z)}\right)^{c(-\alpha^2/2)} .
\]
It is a modular form of weight \(12\) on
\(\mathrm{O}^+(\IIpar)\), where \(\IIpar=\IIlor\oplus U\), and it is the
Weyl--Kac--Borcherds denominator of \(\gFM\).
\end{theorem}

\begin{proof}
Borcherds constructs \(\gFM\) from \(\IIlor\) and proves the displayed
denominator identity. The exponent normalisation follows directly from
\(\eta^{-24}=q^{-1}+24+324q+\cdots\): real roots of norm \(2\) have
exponent \(c(-1)=1\), and null imaginary roots have multiplicity
\(c(0)=24\). Borcherds' singular theta-lift weight formula gives
\[
 \kBKM(\PhiFM)
 =
 \wt(\PhiFM)
 =
 \frac{c(0)}{2}
 =
 \frac{24}{2}
 =
 12 .
\]
\end{proof}

The alternating side has the corresponding Weyl form
\[
 \PhiFM(Z)
 =
 \sum_{w\in W(\gFM)}
 \det(w)\,
 w\!\left(
   e^{2\pi i(\rho,Z)}
   \prod_{n\geq1}
   \left(1-e^{2\pi i(n\rho,Z)}\right)^{24}
 \right),
\]
which is the denominator identity written at the Leech cusp. The input
\(\eta^{-24}\), the lattice \(\IIlor\), the Weyl vector \(\rho\), and the
product above are part of the datum. A comparison with any Calabi--Yau
construction must preserve all four pieces.

\section{The K3 x E Denominator}

The \(K3\times E\) BKM object is governed by the Gritsenko--Nikulin
\(\Delta_5\) denominator. Let
\[
 D_5=64^{-1}\Delta_5 .
\]
In the monic denominator convention,
\[
 \operatorname{den}(\gDelta)=D_5(2Z),\qquad
 \kBKM(\gDelta)=\kBKM(\Delta_5)=5 .
\]
The square normalisations are scalar conventions:
\[
 \Phi_{10}^{\mathrm{un}}=\Delta_5^2,\qquad
 \Phi_{10}^{\mathrm{OP}}=D_5^2=4096^{-1}\Delta_5^2 .
\]
Hence the Oberdieck--Pandharipande scalar branch reads
\[
 Z_{\mathrm{OP/DT}}=-(\Phi_{10}^{\mathrm{OP}})^{-1}
 =
 -D_5^{-2}
 =
 -4096\,\Delta_5^{-2}.
\]
The primitive BKM denominator remains \(D_5(2Z)\).
The Jacobi input for this primitive denominator is
\[
 \phi_{0,1}(\tau,z)=\zeta^{-1}+10+\zeta+O(q).
\]
The doubled input \(2\phi_{0,1}\) has constant term \(20\), hence
Borcherds weight \(10\); it lands on the square line
\(\Phi_{10}^{\mathrm{un}}=\Delta_5^2\). Thus
\(\phi_{0,1}\) supplies the primitive root multiplicities, while
\(2\phi_{0,1}\) supplies the Igusa scalar square.

\begin{proposition}[Primitive CHL constants]
\label{prop:primitive-chl-constants}
For the primitive \(K3\times E\) CHL denominator branch
\[
 N\in\{1,2,3,4,6\},
\]
the Borcherds input constants and BKM weights are
\[
\begin{array}{c|ccccc}
N & 1 & 2 & 3 & 4 & 6\\
\hline
c_N(0) & 10 & 8 & 6 & 4 & 2\\
\kBKM(\Phi_N)=c_N(0)/2 & 5 & 4 & 3 & 2 & 1 .
\end{array}
\]
The \(N=1\) entry is \(\Delta_5\).
\end{proposition}

\begin{proof}
For the \(N=1\) K3 elliptic genus input,
\[
 \phi_{0,1}(\tau,z)=\zeta^{-1}+10+\zeta+O(q),
\]
so \(c_1(0)=10\) in the Eichler--Zagier discriminant-zero convention and
\(\wt(\Delta_5)=5\). For \(N=2,3,4,6\), the Gritsenko--Nikulin CHL
paramodular denominator forms have weights \(4,3,2,1\). Borcherds'
weight formula \(\wt(\Phi_N)=c_N(0)/2\) gives the constants
\((8,6,4,2)\). Thus the full primitive tuple is
\((10,8,6,4,2)\), and the BKM-weight tuple is
\((5,4,3,2,1)\).
\end{proof}

\begin{proposition}[Gritsenko--Clery rows]
\label{prop:gritsenko-clery-row-separation}
The Gritsenko--Clery diagonal-divisor catalogue is indexed by
paramodular triples \((t,N;k)\). The primitive CHL order is a separate
indexing datum.
Its eight rows are
\[
\begin{array}{c|c|c|c}
(t,N;k) & \text{form} & c^{\mathrm{eff}}_{t,N}(0) & \kBKM\\
\hline
(1,1;5) & \Delta_5 & 10 & 5\\
(2,1;2) & \Delta_2 & 4 & 2\\
(1,2;3) & \nabla_3 & 6 & 3\\
(3,1;1) & \Delta_1 & 2 & 1\\
(1,3;2) & \nabla_2 & 4 & 2\\
(4,1;\tfrac12) & \Delta_{1/2} & 1 & \tfrac12\\
(1,4;\tfrac32) & \nabla_{3/2} & 3 & \tfrac32\\
(2,2;1) & Q_1 & 2 & 1 .
\end{array}
\]
Row by row,
\[
 \kBKM(F_{t,N})=\wt(F_{t,N})
 =\frac{c^{\mathrm{eff}}_{t,N}(0)}2 .
\]
The intersection with the primitive \(K3\times E\) CHL ladder is the
\((1,1;5)\) row. The rows \((1,2;3)\) and \((1,3;2)\) have constants
\(6\) and \(4\); they are Gritsenko--Clery paramodular rows, while the
primitive CHL entries at orders \(2\) and \(3\) have constants \(8\) and
\(6\).
\end{proposition}

\begin{proof}
Gritsenko--Clery classify genus-two dd-modular forms with diagonal
divisor by the triple \((t,N;k)\). The weight \(k\) is the Borcherds
weight of the row, so the displayed constants are \(2k\). The primitive
CHL table instead uses the five elliptic-curve-compatible orders
\(\{1,2,3,4,6\}\) and the Gritsenko--Nikulin \(\Delta_{k(N)}^{(N)}\)
tower. Equality of the \((1,1;5)\) row reflects the common
\(\Delta_5\) input; the remaining row labels belong to different
paramodular groups and multiplier systems.
\end{proof}

\begin{proposition}[The \(K3\times E\) invariant split]
\label{prop:k3xe-invariant-split}
For \(Y=K3\times E\), the compact and specialisation invariants are
\[
\begin{array}{c|c|l}
\text{invariant} & \text{value} & \text{source}\\
\hline
\kcat(Y) & 0 &
\chi(\cO_{K3})\chi(\cO_E)=2\cdot0\\
\kch(Y) & 0 &
\sum_q(-1)^q h^{0,q}(Y)=1-1+1-1\\
\kHeis(Y) & 3 &
\text{Heisenberg--Mukai specialisation}\\
\kBKM(\gDelta) & 5 &
c_1(0)/2=10/2\\
\kfib(Y) & 24 &
\rk\,\widetilde H(K3,\Z)=1+22+1 .
\end{array}
\]
These values come from distinct constructions.
\end{proposition}

\begin{proof}
Kunneth multiplicativity gives
\(\kcat(Y)=\chi(\cO_{K3})\chi(\cO_E)=2\cdot0=0\). The compact Hodge
column is \(h^{0,\bullet}(Y)=(1,1,1,1)\), hence
\(\kch(Y)=0\). The Heisenberg--Mukai specialisation contributes the
rank-sum \(2+1=3\). The BKM entry is the Borcherds weight of
\(\Delta_5\). The fibre entry is the Mukai-lattice rank of the K3
factor. The table records five independent constructions.
\end{proof}

\section{K3-Fibre Obstruction}

\begin{proposition}[Single-K3-fibre obstruction]
\label{prop:single-k3-fibre-obstruction}
Let \(X\to B\) be a smooth projective K3-fibred CY\(_3\), and let
\(S\) be a smooth K3 fibre. A specialisation whose real-root lattice is
supplied by the algebraic surface part of one K3 fibre has insufficient
rank for \(\gFM\).
\end{proposition}

\begin{proof}
For a projective K3 surface \(S\),
\[
 \widetilde H_{\mathrm{alg}}(S,\Z)
 =
 H^0(S,\Z)\oplus \operatorname{NS}(S)\oplus H^4(S,\Z),
\qquad
\rk\,\widetilde H_{\mathrm{alg}}(S,\Z)=2+\rho(S),
\]
with \(\rho(S)\leq20\). The \(H^0\oplus H^4\) summand is the Mukai
hyperbolic plane. The algebraic surface part has rank at most \(20\).
The Fake-Monster real-root datum contains the Leech lattice of rank
\(24\) at the cusp \(\rho^\perp/\Z\rho\simeq\Leech\). Thus one K3
surface fibre supplies too small an algebraic surface lattice for the
Leech real-root datum.

The full Mukai lattice \(\widetilde H(S,\Z)\simeq \mathrm{II}_{4,20}\)
also has the wrong signature for the Lorentzian lattice \(\IIlor\) in a
one-fibre construction. The obstruction proved here is the K3-fibre
route obstruction. A global exclusion for all smooth proper CY\(_3\)
categories would require a separate theorem controlling every possible
rank-\(24\) algebraic root lattice in the relevant Hochschild or Hall
charge lattice.
\end{proof}

\section{Dimension-Five Comparison Datum}

Let
\[
 X=K3_1\times K3_2\times E,\qquad
 \Sigma_4=K3_1\times K3_2,\qquad C=E .
\]
Then \(X\) is a compact CY\(_5\) with
\[
 \kcat(X)
 =
 \chi(\cO_{K3_1})\chi(\cO_{K3_2})\chi(\cO_E)
 =
 2\cdot2\cdot0
 =
 0,
\]
and the compact Hodge supertrace also vanishes:
\[
 \kch(X)=\sum_q(-1)^q h^{0,q}(X)
 =(1-1)(1+1)^2=0.
\]

\begin{conjecture}[Conditional \(d=5\) Fake-Monster comparison]
\label{conj:conditional-d5-fake-monster-comparison}
A second-stage specialisation
\[
 \SpCh_{\Sigma_4,E}\bigl(\PhiFA_5(\Perf(X))\bigr)
\]
realising the Fake-Monster denominator branch requires the following
constructed data:
\begin{enumerate}
\item a primitive embedding
\[
 j:\epsilon\IIlor\hookrightarrow
 \widetilde H(K3_1,\Z)\oplus \widetilde H(K3_2,\Z)\oplus U_E,
 \qquad \epsilon\in\{\pm1\},
\]
with \(j(\rho)\) primitive isotropic and
\(j(\rho)^\perp/\Z j(\rho)\simeq \Leech\);
\item a choice of Weyl chamber and positive-root set whose real roots,
imaginary roots, Weyl vector, and multiplicities match
Theorem~\ref{thm:fm-borcherds-denominator-data};
\item a Hall--Drinfeld comparison from the positive-half Hall object of
the CY\(_5\) category to the Borcherds positive half of \(\gFM\);
\item an \(E_5\) factorisation-algebra model and curve specialisation to
an \(E_1\)-chiral algebra on \(E\) whose denominator character is
\(\PhiFM\).
\end{enumerate}
\end{conjecture}

\begin{proof}[Rank evidence]
For a K3 surface,
\(\rk\,\widetilde H(K3,\Z)=24\) and
\(\operatorname{sign}\widetilde H(K3,\Z)=(4,20)\). The doubled Mukai
lattice has rank \(48\) and signature \((8,40)\). With this Mukai sign
convention the embedding test is for \(-\IIlor\), of signature
\((1,25)\); with the opposite Mukai sign it is for \(\IIlor\). Rank
alone leaves room for a sign-compatible Leech summand. The discriminant
form, orientation, and primitive isotropic vector are integral data, and
Nikulin's embedding theorem is the relevant tool once they are
specified.

Rank accommodation is the first numerical test. The comparison
requires the Leech cusp, the Weyl vector, the root multiplicities from
\(\eta^{-24}\), and the automorphic product \(\PhiFM\). It also requires
the Hall--Drinfeld comparison and the factorisation-algebra
specialisation. These are precisely the data listed in the conjecture.
\end{proof}

\begin{remark}[Hodge count on \(K3_1\times K3_2\)]
Kunneth gives
\[
 h^{1,1}(K3_1\times K3_2)
 =
 h^{1,1}(K3_1)h^{0,0}(K3_2)
 +
 h^{0,0}(K3_1)h^{1,1}(K3_2)
 =
 20+20=40 .
\]
The middle Hodge space has
\[
 h^{2,2}(K3_1\times K3_2)
 =
 1+1+1+1+20\cdot20
 =
 404.
\]
For generic projective \(K3_i\) with Picard ranks \(\rho_i\), the
tautological algebraic \(H^{2,2}\)-rank is \(2+\rho_1\rho_2\), with
further Hodge classes on special Noether--Lefschetz loci. A
Leech-cusp comparison is therefore an additional lattice datum beyond
Hodge-number arithmetic.
\end{remark}

\section{Scope Separations}

\begin{proposition}[Four automorphic branches]
\label{prop:four-denominator-branches}
The following branches carry different automorphic data:
\[
\begin{array}{c|c|c|c}
\text{object} & \text{lattice or host} & \text{form} & \text{weight}\\
\hline
\gDelta &
\begin{array}{c}
U\oplus\langle -2\rangle\\
\text{inside }2U\oplus\langle -2\rangle
\end{array}
& D_5(2Z) & 5\\
\Phi_{10}^{\mathrm{OP}} & \text{OP scalar convention} & D_5^2 & 10\\
\gFM & \IIlor=\Leech\oplus U & \PhiFM & 12\\
V^\natural & \text{Monster CFT} & J(\sigma)-J(\tau) & 0 .
\end{array}
\]
The square convention \(\Phi_{10}^{\mathrm{un}}=\Delta_5^2\) relates the
Igusa scalar line to \(\Delta_5\). A comparison between \(\PhiFM\) and
\(\Delta_5\) requires an explicit rational quadratic embedding of type-IV
domains and a divisor-correction calculation for the restricted
Borcherds product.
\end{proposition}

\begin{proof}
The first row is the Gritsenko--Nikulin denominator algebra for the
primitive \(K3\times E\) branch. The second row is a scalar square used
in the OP/DT partition function. The third row is Borcherds'
Fake-Monster denominator from \(\eta^{-24}\) on \(\IIpar\). The fourth
row is the genus-zero Monster denominator. The lattices, Weyl vectors,
input modular forms, and weights differ row by row.
\end{proof}

\begin{proposition}[Chiral and quantum-group scope]
\label{prop:chiral-quantum-group-scope}
For \(d\geq3\), the specialised CY-to-chiral output on a curve is
\(E_1\)-chiral. An \(E_2\) structure belongs to a constructed centre,
Drinfeld double, or module layer. The \(K3\times E\) BKM/Hall--Drinfeld
denominator branch is distinct from the Mukai/Yangian branch. In the
toric local model,
\[
 \CoHA(\C^3)=Y^+(\widehat{\mathfrak{gl}}_1),
\]
the positive half. The full \(\mathcal W_{1+\infty}\) appears on the
double, centre, completion, or evaluation layer.
\end{proposition}

\begin{proof}
The two-stage construction
\[
 \Phi_d=\SpCh_{\Sigma_{d-1},C}\circ\PhiFA_d
\]
specialises a CY\(_d\) factorisation algebra to a chiral algebra on the
curve \(C\). For \(d\geq3\), this curve-level specialised output is
\(E_1\)-chiral. Centre and double constructions are additional algebraic
operations and carry their own hypotheses. The K3 denominator branch is
controlled by Borcherds products and Hall--Drinfeld completion, while
the Mukai/Yangian branch is controlled by stable envelopes, Mukai
lattice symmetries, and current presentations. The \(\C^3\) CoHA
identification is the positive-half shuffle algebra
\(Y^+(\widehat{\mathfrak{gl}}_1)\); passing to
\(\mathcal W_{1+\infty}\) uses the doubled or completed object.
\end{proof}

\section{Proof Obligations}

A constructed dimension-five Fake-Monster comparison must provide the
following mathematical data.
\[
\begin{array}{l|l}
\text{datum} & \text{required content}\\
\hline
\text{Leech-cusp embedding} &
j:\epsilon\IIlor\hookrightarrow
\widetilde H(K3_1)\oplus\widetilde H(K3_2)\oplus U_E\\
\text{Weyl chamber} &
\rho,\ \Delta^+,\ W,\ \operatorname{mult}(\alpha)=c(-\alpha^2/2)\\
\text{Borcherds product} &
\PhiFM=e^{2\pi i(\rho,Z)}
\prod_{\alpha>0}(1-e^{2\pi i(\alpha,Z)})^{c(-\alpha^2/2)}\\
\text{Hall comparison} &
\text{positive half, pairing, coproduct, centre, completion}\\
\text{chiral specialisation} &
\SpCh_{\Sigma_4,E}(\PhiFA_5(\Perf(X)))\text{ as an }E_1\text{-chiral algebra}\\
\text{normalisation} &
\kBKM(\PhiFM)=12,\quad
\kBKM(\Delta_5)=5,\quad
D_5=64^{-1}\Delta_5,\quad
\Phi_{10}^{\mathrm{un}}=\Delta_5^2 .
\end{array}
\]
Until these data are supplied, the proved assertions are Borcherds'
construction of \(\gFM\), the \(\Delta_5\) theorem on \(K3\times E\), and
the single-K3-fibre obstruction above.

\end{document}
